\section{The exact saddle and the product model}\label{ssm:sec:saddle}
For $r\in(0,1)$, let $X_r$ be $\Bin(M-1,r)$ conditioned to be at most $k$, and define
\[
 \mu_M(r)=\E X_r,\qquad
 p(r)=\frac{r^2}{r^2+(1-r)^2}.
\]
We use one saddle for all $f$, chosen by the exact identity
\begin{equation}\label{ssm:eq:exactsaddle}
 \mu_M(r_M)=(M-1)p(r_M),\qquad
 r_M=\frac12-\frac{L_M}{2\sqrt M},\quad 0\le L_M\le1.
\end{equation}

\begin{lemma}\label{ssm:lem:saddle}
For sufficiently large $M$, a root in \eqref{ssm:eq:exactsaddle} exists, uniformly for bounded offsets $k-M/2$. Every such root has $L_M\to\ell$. Moreover,
\begin{equation}\label{ssm:eq:saddlelimits}
 S_M\to\Phi(\ell),\qquad
 \frac{\Var(X_{r_M})}{M}\to v,\qquad
 \frac{\Delta_M}{\sqrt M}\to\frac\ell2,
\end{equation}
where $\Delta_M=(M-1)r_M-\mu_M(r_M)$.
\end{lemma}
\begin{proof}
Write $b=M-1$ and $r=1/2-L/(2\sqrt M)$ with $0\le L\le1$. The binomial central limit theorem is uniform on this compact range of parameters, and its centered moments are uniformly integrable by the usual binomial exponential-moment bound. Its standardized upper endpoint is
\[
 \frac{k-br}{\sqrt{br(1-r)}}=L+O(M^{-1/2}).
\]
Writing $\lambda(L)=\phi(L)/\Phi(L)$, the elementary truncated normal identities are
\[
 \E[Z\mid Z\le L]=-\lambda(L),\qquad
 \Var(Z\mid Z\le L)=1-L\lambda(L)-\lambda(L)^2,
\]
obtained by integration using $\phi'(z)=-z\phi(z)$. Hence
\[
 \mu_M(r)=br-\tfrac12\sqrt M\lambda(L)+o(\sqrt M),
 \qquad p(r)=\tfrac12-LM^{-1/2}+O(M^{-3/2}),
\]
uniformly. Subtraction gives
\begin{equation}\label{ssm:eq:rootlimit}
 \frac{\mu_M(r)-(M-1)p(r)}{\sqrt M}
 =\frac12\{L-\lambda(L)\}+o(1).
\end{equation}
The limiting function is negative at $0$, positive at $1$, and has its unique zero at $\ell$, by \eqref{ssm:eq:ell}. Continuity gives a finite-$M$ root, and compactness plus uniform convergence forces every such root to converge to $\ell$. Finite-$M$ uniqueness is unnecessary.

The binomial cap probability tends to $\Phi(\ell)$. The truncated variance identity, with $\lambda(\ell)=\ell$, gives the variance limit $v=(1-2\ell^2)/4$. Finally the exact saddle yields
\[
 \Delta_M=(M-1)(r_M-p(r_M))=\tfrac12\ell\sqrt M+o(\sqrt M).
\]
All assertions are uniform over bounded offsets by the same compactness argument.
\end{proof}

Set $r=r_M$, $p_0=p(r_M)$, and $m_0=Hp_0$. Under a product law $R_f$, take $M-f$ independent copies of $X_r$ and $f$ independent ordinary $\Bin(M-1,r)$ variables. These are degree coordinates for a comparison calculation, not the degrees of an independent-edge graph. Define
\begin{equation}\label{ssm:eq:Yandshift}
 Y_M=\frac{\sum_iX_i-2m_0}{M},\qquad
 a_M=\E_{R_f}Y_M=\frac{f\Delta_M}{M}.
\end{equation}
The saddle identity ensures that $a_M=0$ for $f=0$.

\begin{lemma}\label{ssm:lem:productlimits}
Along any sequence with bounded $k-M/2$ and $f/\sqrt M\to\tau\in[0,\infty)$,
\begin{equation}\label{ssm:eq:productlimits}
 Y_M\ \Longrightarrow\ \Normal(a,v),\qquad
 a=\tau\ell/2,\qquad
 \gamma(\boldsymbol X)\ \longrightarrow\ v\quad\text{in probability},
\end{equation}
and $R_f(\cC)\to1$. The means $a_M$ are uniformly bounded for $f\le C\sqrt M$, and the exponential-square moment bound \eqref{ssm:eq:shiftUI} applies to $Y_M$ uniformly in that window.
\end{lemma}
\begin{proof}
The total variance divided by $M^2$ tends to $v$: the $M-f$ capped coordinates each contribute $vM+o(M)$, while replacing $f=o(M)$ of them changes the normalized variance by $O(f/M)$. The third absolute centered moments are $O(M^{3/2})$ uniformly by Lemma~\ref{ssm:lem:subgaussian}, so their sum divided by the total variance to the power $3/2$ is $O(M^{-1/2})$. The triangular-array Lyapunov theorem therefore gives a centered normal limit of variance $v$. Equation~\eqref{ssm:eq:saddlelimits} gives $a_M\to a$.

For the sample second moment, center each coordinate at its own mean. Its fourth moment is $O(M^2)$, so the variance of $M^{-2}\sum_i(X_i-\E X_i)^2$ is $O(M^{-1})$. Its mean tends to $v$. The difference between the two group means is $\Delta_M=O(\sqrt M)$, whose additional squared contribution divided by $M^2$ is $O(f\Delta_M^2/M^2)=O(f/M)=o(1)$. Cross terms vanish by Cauchy--Schwarz. Subtracting the square of the sample mean error changes the normalized second moment by $o_{\PP}(1)$, and replacing $M^2$ by $(M-1)^2$ has no effect. This proves the assertion about $\gamma$.

Every coordinate's mean differs from $M/2$ by $O(\sqrt M)$. The single-coordinate bound \eqref{ssm:eq:mgf} therefore gives
\[
 R_f(\cC^c)\le2M\exp\{-cM^{4\delta}\}=o(1)
\]
for large $M$. The square-moment statement is Corollary~\ref{ssm:cor:squareUI}, since $Y_M=Z_M+a_M$ with $a_M$ bounded.
\end{proof}
